% \section{Dataset annotation}
% \label{app:dataset_annotation}
% \subsection{Human annotators}
% \label{app:human_annotators}
% We introduce human annotators in \ours data collection.
% For the StackExchange datasets,
% Hongjin Su annotated/adapted \econ, \sustainable, \psychology, \robotics and \pony; Han-yu Wang annotated \biology; Haisu Liu annotated \stackoverflow; Qilin Liao annotated \earth.
% In addition to the author review, we also invite domain experts (PhD students) to review the data: Yun Han reviews \sustainable, Xiaoru Teng reviews \psychology, Cong Gao reviews \econ, Shengyu Wang, Xiaodong Wei and Yan Pan review \biology.

% For the TheoremQA dataset, three of the authors--Quan Shi, Zachary S. Siegel, and Michael Tang--manually checked the rewritten questions for consistency and solvability.
% For the \aops dataset, Howard Yen, Quan Shi, Zachary S. Siegel, and Michael Tang browsed through the \aops Wiki and collected the topics (i.e., the theorems and problem-solving techniques) and their respective problem statements and solutions. 
% For the \leetcode dataset, Howard Yen labeled questions on whether they are grounded in real-world concepts for checking human-model agreement. More details are described in \S\ref{app:data:leetcode}.
% All involved authors are undergraduate or graduate students in computer science. 

\subsection{LLM usage}
\label{app:llm_usage}
There is no specific procedure for the annotators to use LLMs. 
The LLMs serve a tool to help annotators understand queries and documents, i.e., whenever they fail to understand something, they ask LLMs for clarification, explanation, etc.
They are also used to summarize the content of StackExchange posts for searching negative documents in Google for StackExchange. 
To diversify the search results, the annotator prompts LLMs with role-playing scenarios (e.g., as a biology student) to generate keywords of the posts. 
% These keywords are then used for Google searches to collect additional web pages.
% into a few keywords, which are then used to search Google for documents with lexical similarity or semantic similarity, but irrelevant to the post.
The responses from LLMs are not included in any of \ours datasets.
Since the responses from LLMs are not guaranteed to be correct, the annotators always search for trustworthy sources to verify the information.
The LLMs used in the annotation include ChatGPT\footnote{https://chatgpt.com/}, GPT-4, and Claude-3.

\subsection{Sensitive information}
\label{app:sensitive}
All the data in \ours are manually collected, carefully verified, and reviewed to remove any personally identifiable information or offensive content.